\section*{При каких температурах измерены обе ветви}
В таблице~1 статьи Кайнумы, Пальма и Индена приведены составы на
Ti-богатой границе между разупорядоченным ГПУ-раствором \(\alpha\) и
упорядоченной фазой \(\alpha_2\) типа D0$_{19}$.
При 1000 и 1050~\(^\circ\)C авторы отжигали по 200~ч диффузионные пары
с исходными составами 15 и 25~ат.\% Al. Равновесные составы определяли
по концентрационным профилям у границы фаз электронно-зондовым
микроанализом. При 1100~\(^\circ\)C приведены точечные измерения в
двухфазном образце. Знак \(\pm\) в таблице передаёт указанную в
первоисточнике точность измерений.
\href{https://doi.org/10.1016/0966-9795(94)90018-3}{Kainuma, Palm, Inden (1994), табл.~1, с.~326.}

\begin{center}\small
\begin{tabular}{@{}rrrrrr@{}}
\toprule
&&\multicolumn{2}{c}{Эксперимент, ат.\% Al}&
\multicolumn{2}{c}{Liang, ат.\% Al}\\
\cmidrule(lr){3-4}\cmidrule(l){5-6}
\(T,\ ^\circ\)C & \(T,\) К & \(\alpha\) & \(\alpha_2\) & \(\alpha\) & \(\alpha_2\)\\
\midrule
1000 & 1273,15 & \(19{,}5\pm0{,}5\) & \(23{,}0\pm0{,}3\) & 19,35 & 23,19\\
1050 & 1323,15 & \(21{,}3\pm0{,}3\) & \(23{,}5\pm0{,}3\) & 20,74 & 23,96\\
1100 & 1373,15 & \(23{,}5\pm0{,}2\) & \(24{,}9\pm0{,}2\) & 22,64 & 25,06\\
1150 & 1423,15 & \(24{,}7\pm0{,}2\) & \(25{,}7\pm0{,}2\) & 26,09 & 27,22\\
\bottomrule
\end{tabular}
\end{center}
Две ближайшие к 1300~К строки соответствуют разным изотермам. При
1273,15~К эксперимент разрешает интервал шириной \(3{,}5\)~п.п., при
1323,15~К его ширина составляет \(2{,}2\)~п.п. Поэтому для прямого
атомистического теста выбрана первая температура. В таблице~1 независимой
работы Онумы и соавт. при 1373,15~К для \(\alpha/\alpha_2\) приведены
23,5 и 24,8~ат.\% Al. Их строка при 1273,15~К относится к
\(\alpha_2/\gamma\), то есть к границе с другой фазой.
\href{https://doi.org/10.1016/S1359-6454(00)00118-X}{Ohnuma et al. (2000), табл.~1.}

Для 600~К прямой пары составов в проверенных экспериментальных
таблицах и нанесённых в проекте точках Блэкберна не найдено.
Для этой изотермы количественный внешний ориентир даёт база Liang.
Во всех строках CALPHAD рассчитан заново при \(p=0\), с базой Liang
и десятью фазами бинарного сечения в \texttt{pycalphad}~0.11.2.
Отбиралась Ti-богатая ветвь с \(x_\alpha<x_{\alpha_2}\).
Для 600 и 1273,15~К расчёт только с \texttt{HCP\_A3} и
\texttt{ALTI3\_D019} дал те же контакты, что и расчёт с десятью фазами.
База и обработка фиксируют, с какой именно термодинамической оценкой
сопоставляется MTP. \href{https://doi.org/10.1007/s11669-024-01124-9}{Liang et al. (2024).}
